\documentclass[tikz,border=8pt]{standalone}
\usepackage{tikz}
\usepackage{amsmath}
\usepackage{amssymb}
\usepackage{pgfplots}
\pgfplotsset{compat=1.18}
\usepackage{ctex}
\usetikzlibrary{calc,positioning,arrows.meta,matrix,trees}

% LC 191 / 338 位计数 Brian Kernighan 算法
% n=11=1011，每步 n &= (n-1) 消除最低位1，统计步数

\begin{document}
\begin{tikzpicture}[
  every node/.style={font=\small},
  bitcell/.style={draw=gray!50, rectangle, minimum width=0.52cm, minimum height=0.55cm, thin, fill=white},
  onebit/.style={draw=blue!60, rectangle, minimum width=0.52cm, minimum height=0.55cm, thin, fill=blue!15},
  rembit/.style={draw=red!60, rectangle, minimum width=0.52cm, minimum height=0.55cm, thin, fill=red!15},
  cell/.style={draw=gray!60, rectangle, minimum width=1.8cm, minimum height=0.65cm, thick, fill=white},
  hcell/.style={draw=orange!70, rectangle, minimum width=1.8cm, minimum height=0.65cm, thick, fill=orange!10},
  gcell/.style={draw=green!60!black, rectangle, minimum width=1.8cm, minimum height=0.65cm, thick, fill=green!10},
  ptr/.style={->,>=Stealth,thick},
]

\def\bw{0.55}  % bit cell width spacing

%% 标题
\node[font=\large\bfseries] at (5.5, 1.5)
  {位计数 Brian Kernighan 算法（LC 191/338）};
\node[font=\small,gray] at (5.5, 0.9)
  {n=11=1011$_2$，每步 n \&= (n-1) 消除最低位的 1，统计步数即 popcount};

%% 表头
\node[font=\scriptsize\bfseries,gray] at (1.5, 0.2) {n（十进制）};
\node[font=\scriptsize\bfseries,gray] at (3.8, 0.2) {n-1（十进制）};
\node[font=\scriptsize\bfseries,gray] at (6.2, 0.2) {n 二进制（4位）};
\node[font=\scriptsize\bfseries,gray] at (8.8, 0.2) {n\&(n-1) 结果};
\node[font=\scriptsize\bfseries,gray] at (10.8, 0.2) {count};

\draw[gray!40] (0.5, -0.0) -- (12.0, -0.0);

%% 步骤行1：n=11 (1011)，n-1=10 (1010)，n&(n-1)=10 (1010)，消最低位1
\def\rowA{-0.8}
\node[font=\scriptsize] at (1.5, \rowA) {n = 11};
\node[font=\scriptsize] at (3.8, \rowA) {n-1 = 10};
% n=11: 1011
\node[onebit] at (5.3+0*\bw, \rowA) {1};
\node[onebit] at (5.3+1*\bw, \rowA) {0};
\node[onebit] at (5.3+2*\bw, \rowA) {1};
\node[rembit] at (5.3+3*\bw, \rowA) {1};
\node[font=\tiny,red!70,anchor=north] at (5.3+3*\bw, \rowA-0.3) {消除};
% n&(n-1) = 10 = 1010
\node[font=\scriptsize] at (8.8, \rowA) {1010 = 10};
\node[gcell,minimum width=1.0cm] at (10.8, \rowA) {1};

%% 步骤行2：n=10 (1010)，n-1=9 (1001)，n&(n-1)=8 (1000)
\def\rowB{-1.9}
\node[font=\scriptsize] at (1.5, \rowB) {n = 10};
\node[font=\scriptsize] at (3.8, \rowB) {n-1 = 9};
% n=10: 1010
\node[onebit] at (5.3+0*\bw, \rowB) {1};
\node[bitcell] at (5.3+1*\bw, \rowB) {0};
\node[rembit] at (5.3+2*\bw, \rowB) {1};
\node[bitcell] at (5.3+3*\bw, \rowB) {0};
\node[font=\tiny,red!70,anchor=north] at (5.3+2*\bw, \rowB-0.3) {消除};
% n&(n-1) = 8 = 1000
\node[font=\scriptsize] at (8.8, \rowB) {1000 = 8};
\node[gcell,minimum width=1.0cm] at (10.8, \rowB) {2};

%% 步骤行3：n=8 (1000)，n-1=7 (0111)，n&(n-1)=0
\def\rowC{-3.0}
\node[font=\scriptsize] at (1.5, \rowC) {n = 8};
\node[font=\scriptsize] at (3.8, \rowC) {n-1 = 7};
% n=8: 1000
\node[rembit] at (5.3+0*\bw, \rowC) {1};
\node[bitcell] at (5.3+1*\bw, \rowC) {0};
\node[bitcell] at (5.3+2*\bw, \rowC) {0};
\node[bitcell] at (5.3+3*\bw, \rowC) {0};
\node[font=\tiny,red!70,anchor=north] at (5.3+0*\bw, \rowC-0.3) {消除};
% n&(n-1) = 0
\node[font=\scriptsize] at (8.8, \rowC) {0000 = 0};
\node[gcell,minimum width=1.0cm] at (10.8, \rowC) {3};

%% 步骤行4：n=0，结束
\def\rowD{-4.1}
\node[font=\scriptsize] at (1.5, \rowD) {n = 0};
\node[font=\scriptsize,gray] at (3.8, \rowD) {（终止）};
\node[bitcell] at (5.3+0*\bw, \rowD) {0};
\node[bitcell] at (5.3+1*\bw, \rowD) {0};
\node[bitcell] at (5.3+2*\bw, \rowD) {0};
\node[bitcell] at (5.3+3*\bw, \rowD) {0};
\node[font=\scriptsize,gray] at (8.8, \rowD) {n==0 停止};
\node[draw=green!60!black,fill=green!15,rounded corners=2pt,
      font=\scriptsize\bfseries,inner sep=2pt] at (10.8, \rowD) {3};

\draw[gray!40] (0.5, -4.55) -- (12.0, -4.55);

%% 位列标注
\node[font=\tiny,gray] at (5.3+0*\bw, -5.0) {bit3};
\node[font=\tiny,gray] at (5.3+1*\bw, -5.0) {bit2};
\node[font=\tiny,gray] at (5.3+2*\bw, -5.0) {bit1};
\node[font=\tiny,gray] at (5.3+3*\bw, -5.0) {bit0};

%% 图例
\node[onebit,minimum width=0.5cm] at (0.9, -5.3) {1};
\node[font=\scriptsize,anchor=west] at (1.2, -5.3) {值为1的位};
\node[rembit,minimum width=0.5cm] at (0.9, -5.9) {1};
\node[font=\scriptsize,anchor=west] at (1.2, -5.9) {被消除的最低位1};

%% n &(n-1) 原理说明
\node[draw=gray!40,fill=white,thin,rounded corners=2pt,
      font=\scriptsize,inner sep=6pt,align=left]
  at (5.5, -6.9)
  {\textbf{n \& (n-1) 的效果：}\\[3pt]
   n     = \ldots X 1 0 0 0（最低位1位于某位置）\\[2pt]
   n-1   = \ldots X 0 1 1 1（该位变0，低位全变1）\\[2pt]
   n\&(n-1) = \ldots X 0 0 0 0（消除最低位的1）};

%% 结果与算法说明
\node[draw=green!60!black,fill=green!8,very thick,rounded corners=4pt,
      font=\large\bfseries]
  at (5.5, -8.4) {popcount(11) = 3（二进制 1011 含 3 个 1）};

\node[draw=gray!50,fill=white,rounded corners=3pt,font=\small,
      inner sep=8pt,align=left]
  at (5.5, -9.9)
  {\textbf{Brian Kernighan 位计数 O(k) 时间（k=1的个数）：}\\[4pt]
   \texttt{count=0; while n!=0: count++; n \&= (n-1)}\\[2pt]
   每次迭代消除 n 中最低位的1，迭代次数=1的个数};

\end{tikzpicture}
\end{document}
